\begin{lemma}
Let $n\ge 3$, and let $A=(a_{i,j})_{i\in[3],j\in[n]}$ be a $3\times n$
matrix with entries in $[0,1]$.  Put $c_j=\sum_{i=1}^3 a_{i,j}$.  Suppose
$c_j\le 1$ for every $j$, $\sum_{j=1}^n c_j\le 3$, and
$1\ge c_1\ge c_2\ge\cdots\ge c_n$.  Then
\[
\frac{1}{9}\sum_{j=1}^{n}\sum_{1\le i<i'\le3}a_{i,j}a_{i',j}
+\frac{1}{27}\operatorname{perm}(A_{[3]})
-\frac{1}{27}\sum_{i=1}^3\sum_{j=1}^3a_{i,j}
\le \frac{2}{3^5},
\]
where $A_{[3]}$ is the leading $3\times3$ submatrix.
\end{lemma}
\begin{proof}
Let
\[
Q=\sum_{j=1}^{n}\sum_{1\le i<i'\le3}a_{i,j}a_{i',j},\qquad
P=\operatorname{perm}(A_{[3]}),\qquad
S=c_1+c_2+c_3 .
\]
Multiplying the desired inequality by $27$, it is enough to prove
$3Q+P-S\le 2/9$.

The leading-column estimate is supplied by
\noderef{ThreeByThreeCenteredColumnBound}, applied to $A_{[3]}$;
its hypotheses hold by nonnegativity and the three column-sum bounds.
We recall its centered-column calculation. Write the first three columns in the form
\[
a_{\cdot,1}=\frac{c_1}{3}\mathbf 1+u,\qquad
a_{\cdot,2}=\frac{c_2}{3}\mathbf 1+v,\qquad
a_{\cdot,3}=\frac{c_3}{3}\mathbf 1+w,
\]
where the coordinates of each of $u,v,w\in\mathbb R^3$ sum to zero.  Put
$R=\|u\|_2^2+\|v\|_2^2+\|w\|_2^2$.  If $Q_1$ is the contribution to $Q$ from
the first three columns, then, for a vector $z$ with coordinate sum $c$,
\[
\sum_{i<i'}z_iz_{i'}=\frac12\left(c^2-\sum_i z_i^2\right).
\]
Applying this identity to the three displayed columns gives
\[
3Q_1=c_1^2+c_2^2+c_3^2-\frac32R .
\]

Expanding the permanent of the leading $3\times3$ matrix around the three
constant columns gives
\[
P=\frac29c_1c_2c_3
-\frac{c_3}{3}\langle u,v\rangle
-\frac{c_2}{3}\langle u,w\rangle
-\frac{c_1}{3}\langle v,w\rangle
+2\sum_{i=1}^3u_iv_iw_i .
\]
Since $c_1,c_2,c_3\le1$, Cauchy--Schwarz bounds the absolute values of the
three bilinear terms by at most $R/3$ in total.  Also
$0\le a_{i,j}\le c_j\le1$, so every coordinate of $u,v,w$ has absolute value
at most $2/3$.  Therefore
\[
2\left|\sum_i u_iv_iw_i\right|
\le \frac43\sum_i|u_iv_i|
\le \frac23(\|u\|_2^2+\|v\|_2^2)
\le \frac23R .
\]
Thus
\[
P\le \frac29c_1c_2c_3+R.
\]
Combining this with the formula for $3Q_1$ and discarding the nonpositive term
$-R/2$ yields
\[
3Q_1+P\le c_1^2+c_2^2+c_3^2+\frac29c_1c_2c_3 .
\]

For each $j\ge4$,
\[
3\sum_{i<i'}a_{i,j}a_{i',j}\le c_j^2,
\]
because
$\sum_{i<i'}x_ix_{i'}\le \frac13(x_1+x_2+x_3)^2$ for nonnegative
$x_1,x_2,x_3$.  Since the columns are ordered, $c_j\le c_3$ for $j\ge4$; and
since the total column sum is at most $3$,
\[
\sum_{j=4}^n c_j^2\le c_3\sum_{j=4}^n c_j
\le c_3(3-c_1-c_2-c_3).
\]
Writing $x=c_1$, $y=c_2$, and $z=c_3$, we get
\[
3Q+P-S
\le x(x-1)+y(y-1)+z\left(2-x-y+\frac29xy\right).
\]
The coefficient of $z$ is nonnegative: from $0\le x,y\le1$,
$2-x-y\ge0$, and the remaining term is nonnegative.  Because $z\le y$, the
last display is at most
\[
x(x-1)+y(y-1)+y\left(2-x-y+\frac29xy\right)
=x(x-1)+y\left(1-x+\frac29xy\right).
\]
Its value at $y=x$ minus its value at $y$ equals
\[
(x-y)\left(1-x+\frac29x(x+y)\right)\ge0,
\]
since $0\le y\le x\le1$. Hence it is at most its value at $y=x$:
\[
x(x-1)+x\left(1-x+\frac29x^2\right)=\frac29x^3\le\frac29 .
\]
This proves $3Q+P-S\le2/9$, and therefore the stated inequality.
\end{proof}
